\documentclass{article}
\usepackage{amsmath, amssymb, amsthm}
\usepackage{hyperref}

\title{CNSH Runtime Governance Theory}
\author{UID9622}
\date{\today}

\begin{document}
\maketitle

\begin{abstract}
(请填写摘要)
\end{abstract}

\newtheorem{definition}{Definition}
\newtheorem{axiom}{Axiom}
\newtheorem{lemma}{Lemma}
\newtheorem{theorem}{Theorem}
\newtheorem{corollary}{Corollary}

\begin{axiom}[Axiom 1]
No execution without audit
\begin{equation}
No execution without audit
\end{equation}
\end{axiom}

\begin{definition}[Definition 1]
Digital Root is the modulo-9 reduction of any integer
\begin{equation}
dr(n)=1+((n-1) \bmod 9)
\end{equation}
\end{definition}

\begin{lemma}[Lemma 1]
If dr(n) ∈ {1,2,4,5,7,8}, then n is in GREEN state
\begin{equation}
If dr(n) ∈ {1,2,4,5,7,8}, then n is in GREEN state
\end{equation}
\begin{proof}
By definition of tri-color governance set G
\end{proof}
\end{lemma}

\begin{theorem}[Theorem 1]
Every CNSH runtime execution is traceable via DNA chain
\begin{equation}
Every CNSH runtime execution is traceable via DNA chain
\end{equation}
\begin{proof}
By induction on the DNAStack push chain
\end{proof}
\end{theorem}

\begin{corollary}[Corollary 1]
All GREEN state executions pass without manual review
\begin{equation}
All GREEN state executions pass without manual review
\end{equation}
\end{corollary}

\end{document}